\chapter{Operational semantics and expression traces}
A program adds the integers from $N$ down to 1. For $N=3$, its answer should
be 6. But suppose an optimizer decrements the counter before adding it.
The program is still well typed, every instruction is legal, and the answer
becomes 3. Types alone did not preserve the computation.

Chapter 14 defined typed instructions and snapshot-based state ownership.
Here we expose execution one transition at a time. We use scalar integer
registers to make control and effects visible without tensor-library details.
This is a new, explicitly sequential teaching machine, not a silent change
to Chapter 14's snapshot semantics. Its instruction alphabet is deliberately
small enough to derive, implement and exhaustively trace.

\section{Specify configurations before instructions}
\index{configuration}\index{operational semantics}
A program $P$ is an immutable tuple of instructions. A configuration
\begin{equation}
 C=(p,\rho)
\end{equation}
contains a program counter $p$ and an environment $\rho$ mapping register names
to mathematical integers. The program is fixed during a run. The counter
selects the next instruction, not the number of instructions already executed.
An event records the instruction and its before/after configurations.

\Listing{Configuration}
The helper \texttt{configuration} validates integer values and stores register
pairs in sorted order, making snapshots immutable and comparisons deterministic.
Python integers avoid fixed-width overflow in this reference; their arithmetic
and storage are not constant cost for arbitrarily large values.
\Plate{01-machine}{A configuration points to one instruction and stores register
values. The step relation reads both and returns a new configuration plus an
event. No biological process is implied by the term expression.}{fig:machine}

We use five instruction forms:
\begin{center}
\begin{tabular}{ll}\toprule
Instruction & Effect\\\midrule
$\operatorname{add}(a,b)$ & $\rho(a)\gets\rho(a)+\rho(b)$; advance\\
$\operatorname{dec}(a)$ & $\rho(a)\gets\rho(a)-1$; advance\\
$\operatorname{jump}(j)$ & set counter to $j$\\
$\operatorname{jz}(a,j)$ & branch to $j$ if $\rho(a)=0$, otherwise advance\\
$\operatorname{halt}$ & return a halted outcome\\\bottomrule
\end{tabular}
\end{center}
The complete program is validated before execution: every register exists,
each instruction has the right arity, and every branch target is a valid index.
This prevents a disabled or unvisited branch from hiding malformed code.
\Listing{validate}
Validation is not termination analysis. A legal self-jump can execute forever,
and a legal final non-halt instruction can fall off the program.

\section{Derive the small-step relation}
\index{small-step relation}\index{effects}
Write $C\xrightarrow{e}_P C'$ when executing the selected instruction produces
event $e$ and next configuration $C'$. For an addition,
\begin{equation}
 \frac{P[p]=\operatorname{add}(a,b)}
 {(p,\rho)\longrightarrow
 (p+1,\rho[a\mapsto\rho(a)+\rho(b)])}.
 \label{eq:addstep}
\end{equation}
All right-hand-side values are read from the old environment; if $a=b$, this
doubles the old value. Registers not named by the update remain unchanged.
For zero-branch,
\begin{equation}
 p'=\begin{cases}j&\rho(a)=0,\\p+1&\rho(a)\ne0.\end{cases}
 \label{eq:branch}
\end{equation}
For decrement, the new value is $\rho(a)-1$, including negative results.
Halting is a terminal action, not ordinary fallthrough. The implementation
uses counter $|P|$ as a halted checkpoint, together with a separate status.

\Listing{step}

There is at most one applicable rule for a validated instruction and environment.
Thus the machine is deterministic: the same program and configuration produce
the same next configuration and event. The argument depends on exact integer
operations and a fixed program. Random routing, asynchronous external effects
or floating-point nondeterminism would require additional semantics.

The public runner validates inputs; \texttt{step} is an internal primitive
assuming that validation. It cannot safely execute arbitrary unvalidated tuple
structures merely because its body is short.

\section{A loop invariant proves the result; a ranking function proves termination}
\index{loop invariant}\index{ranking function}
The program has indices 0 through 4:
\[
\begin{array}{rll}
0&\operatorname{jz}(n,4)&\text{check completion}\\
1&\operatorname{add}(s,n)&\text{accumulate current counter}\\
2&\operatorname{dec}(n)&\text{decrease counter}\\
3&\operatorname{jump}(0)&\text{return to check}\\
4&\operatorname{halt}&\text{finish}.
\end{array}
\]
Initialize $n=N\geq0$, $s=0$. At each visit to index 0, the invariant is
\begin{equation}
 s+\frac{n(n+1)}2=\frac{N(N+1)}2,\qquad 0\leq n\leq N.
 \label{eq:sum}
\end{equation}
Initially it is immediate. One body iteration replaces $(s,n)$ by
$(s+n,n-1)$. Substitution gives
\[
 s+n+\frac{(n-1)n}{2}=s+\frac{n(n+1)}2,
\]
so the invariant is preserved at the next loop head.
When $n=0$, it implies $s=N(N+1)/2$.

\Plate{02-loop}{The guarded loop checks zero before changing state.
The rank $n$ decreases once per body traversal. It need not decrease on the
guard or back edge individually.}{fig:loop}

Correctness on termination does not itself prove termination. At loop heads
with $n>0$, rank $n$ is a nonnegative integer that decreases strictly after a
body. It therefore reaches zero after exactly $N$ bodies.
Each body cycle executes four instructions, followed by a final zero check
and halt, so
\begin{equation}
 T(N)=4N+2.
 \label{eq:steps}
\end{equation}
If $N<0$, decrement moves away from zero. The program remains syntactically
valid but the precondition for this termination proof fails.
The runner accepts integer states generally; the summation contract requires
the caller to enforce $N\geq0$.

\begin{center}\input{\Generated/trace.tex}\end{center}
This generated trace for $N=3$ reaches $s=6$ in fourteen transitions.
An event is recorded for the halt instruction too. Changing that convention
would change the exact cost formula.

\section{Fuel bounds observation, not the program's meaning}
\index{fuel}\index{termination}
\Listing{run}
The runner returns one of three statuses. \texttt{halt} means an explicit halt
was executed. \texttt{exhausted} means the provided transition budget was used
without observing halt. \texttt{error} means no valid next instruction exists
after a step, such as falling beyond the last non-halt instruction.
Malformed programs are rejected before running, rather than being classified
as a runtime outcome.

\Plate{05-outcomes}{A bounded run separates an explicit halt, an invalid next
step and budget exhaustion. Exhaustion preserves a checkpoint that can be
resumed; it is not a proof of divergence.}{fig:outcomes}

For a valid initial configuration, zero fuel returns an empty trace and
exhaustion. Fuel $4N+1$ stops just before the halt instruction in the summation
example; $4N+2$ observes it. A self-jump also exhausts any finite fuel supply,
but these equal statuses do not imply equal termination behavior.

The bound controls transitions, not memory or wall-clock time. Integer sizes,
event storage and validation cost are outside its count. The implementation
is not a sandbox for executing untrusted programs; it interprets a fixed,
small instruction alphabet and performs no external I/O.

A resumable checkpoint contains the counter and complete register environment.
Run seven steps at $N=4$, then resume the checkpoint for eleven more:
the concatenated event sequence equals the eighteen-step uninterrupted run.
This follows by determinism and can be tested at every split before halt.
A checkpoint after halt is terminal; the public API does not restart it as
an ordinary instruction.

\section{Well-typed transformations can violate effects}
\index{read-write dependency}\index{optimization legality}
Compare $\operatorname{add}(s,n);\operatorname{dec}(n)$ with the reverse order.
At $(s,n)=(0,3)$ they yield $(3,2)$ and $(2,2)$.
Both preserve integer types. They differ because addition reads $n$ while
decrement writes $n$.
\Plate{03-effects}{Reordering a read across a write changes the result.
The dependency arrow is an ordering constraint, not a data tensor.
Equal types and equal instruction counts do not establish equivalence.}{fig:effects}

For two pure deterministic state updates $A,B$ with read sets $R_A,R_B$
and write sets $W_A,W_B$, a sufficient noninterference condition is
\begin{equation}
 W_A\cap(R_B\cup W_B)=\varnothing,\qquad
 W_B\cap(R_A\cup W_A)=\varnothing.
 \label{eq:commute}
\end{equation}
Under these assumptions neither update changes values the other observes
or writes, so the two orders commute. This is sufficient, not necessary:
some overlapping arithmetic updates commute for algebraic reasons.

Control effects and observability also matter. Exchanging two independent
register updates can preserve the final environment while changing a full
instruction trace. If the specification exposes every event, trace equality
is stronger than final-value equality. A compiler must declare which
observations it preserves: results, exceptions, I/O, state, or a specified
projection of a trace.

Modern MLIR structured-control-flow documentation makes region arguments,
conditions and yielded values explicit \cite{mlir}. StableHLO likewise states
dataflow execution ordering and uses token dependencies for side effects
\cite{stablehlo}. These are relevant production design patterns, not evidence
that this register machine implements either representation.

\section{Replay checks a claim rather than trusting a log}
\index{trace replay}
An event stores before counter, instruction, before registers, after counter
and after registers. Semantic replay recomputes the expected event from the
current configuration, compares the entire record, and advances only on equality.
\Listing{replay}
\Plate{04-replay}{A replay checker recomputes each transition from trusted program
and initial state. It compares the claimed event with the expected event.
A plausible final answer alone cannot certify the intervening trace.}{fig:replay}

Duplicating or changing an event is detected when the next event no longer
matches the derived transition. A valid prefix is still a valid partial trace:
replay does not claim completion unless the caller separately checks that its
last event is halt and its final configuration agrees with the reported result.
Removing an entire suffix can therefore preserve replay validity while destroying
a completeness claim.

The trace is not cryptographically authenticated. An adversary who changes
the program, initial state and all events consistently can produce another valid
trace. Retain content identities and an external trust mechanism when provenance
matters. A textual version label does not supply either automatically.

Tests compare the summation result with $N(N+1)/2$ for $N=0,\ldots,50$,
check exact fuel boundaries, resume a split run, replay its events and reject
a tampered record. Negative controls include a self-loop, a reordered body,
undefined registers, invalid jumps and fallthrough without halt.

\section{Exercises with worked answers}
\begin{enumerate}
\item \textbf{Predict.} Run $N=0$. \emph{Answer:} Zero check and halt execute;
$s=0$, with two transitions.
\item \textbf{Derive.} Prove Equation~\ref{eq:sum}. \emph{Answer:} The body adds $n$
and replaces its triangular remainder by $(n-1)n/2$; the total is unchanged.
\item \textbf{Terminate.} Why is $n$ a valid rank only at loop heads?
\emph{Answer:} It decreases across each complete body, but not on every control step.
\item \textbf{Count.} Give exact fuel for $N=10$. \emph{Answer:} $4(10)+2=42$.
\item \textbf{Distinguish.} Does exhausted fuel prove an infinite loop?
\emph{Answer:} No; the summation run with one step too little fuel is a counterexample.
\item \textbf{Precondition.} What fails for $N=-1$?
\emph{Answer:} The nonnegative rank premise; successive decrements never reach zero.
\item \textbf{Effects.} Reorder the loop's add and decrement.
\emph{Answer:} The new sum is $N(N-1)/2$, although the final counter is still zero.
\item \textbf{Aliasing.} Interpret $\operatorname{add}(s,s)$.
\emph{Answer:} Read the old value twice and store $2s$, not a sequentially updated
right-hand side.
\item \textbf{Replay.} Can a valid prefix establish completion?
\emph{Answer:} No; require a final halt and agreement with the claimed final state.
\item \textbf{Implement.} Test every split of a finite run.
\emph{Rubric:} Resume counter and complete environment, compare concatenated events,
handle zero-length prefixes, and treat a halted checkpoint as terminal.
\item \textbf{Optimize.} State what observations a scheduler preserves.
\emph{Rubric:} Specify final values, effects and trace projection; test a read/write
conflict and an independent pair. Do not equate output parity with full trace parity.
\item \textbf{Budget.} Why is fuel not a memory sandbox?
\emph{Answer:} Integer payloads and retained events can grow; one logical step
does not impose a byte or time limit.
\end{enumerate}
\section{Terms and the next contract}
A \Term{configuration} is the machine's current counter and environment.
A \Term{small step} is one rule application. A \Term{ranking function} proves
progress toward termination. \Term{Observational equivalence} depends on which
outcomes are visible. A \Term{checkpoint} retains enough state to resume.
README commands reproduce the reference and its tables.
The next chapter attaches resource ledgers to executions without confusing
a logical step with arithmetic work, peak memory or measured latency.
